\begin{frame}{FAQ (10.3)}
  \small
  \begin{itemize}
    \item \textbf{분할이 분류보다 더 어려운 이유?} 격자 \textbf{위치}
      는 고정되어 NMS 문제가 없지만 (1) \textbf{라벨 비용}:
      분류는 이미지 1장에 라벨 1개, 분할은 \textbf{픽셀마다}
      라벨(의료 영상 1장에 수백만 개, 전문의가 직접 그려야 함).
      (2) \textbf{경계}: 7$\times$7 격자(풀링 5번) 해상도로는
      ``도로의 가장자리''가 흐려진다 $\to$ U-Net이 스킵으로
      해상도를 \textbf{되찾아야} 함. (3) \textbf{클래스
      불균형} 극단: 자율주행 프레임에서 ``보행자'' 픽셀은
      0.1\% 미만 --- 정확도 99\% 모델이 보행자를 \textbf{전부}
      ``도로''로 분류할 수도 있음(Week 2의 정확도 함정,
      픽셀 스케일에서 재발).
    \item \textbf{미세조정 학습률이 낮은 이유?} (실수 2) ---
      사전학습 가중치는 좋은 손실 지역 위에 있고, 큰 학습률은
      그 지역을 \textbf{벗어나} 범용 특징을 파괴. 작은
      걸음(1/10$\sim$1/100)은 ``\textbf{좋은 점의 근처에서}''
      새 작업에 적응하는 것이지, 처음부터 배우는 게 아님.
    \item \textbf{데이터 100장뿐인데 전이학습을 해도 되나?}
      그것이 전이학습의 \textbf{주식 용례}. 전제는 새 도메인이
      사전학습 도메인과 ``너무 다르지 않다''(가정 풍경
      $\to$ 주차장 풍경은 OK, RGB 자연영상 $\to$ X-ray는
      큰 점프). 100장 수준에서는 \textbf{특징 추출}(헤드만
      학습)이 가장 안정적.
    \item \textbf{구별할 클래스가 ImageNet에 없는데?} \textbf{되
     며, 그게 표준 경우}. 앞쪽 층은 ``고양이를 구분하는
      능력''이 아니라 \textbf{``모서리를 보는 능력''}이라,
      ImageNet 1,000개 클래스에 없던 사물도 그 사물의
      \textbf{모서리·질감}으로 앞쪽 층이 특징을 뽑아냄.
      바꿔야 할 것은 \textbf{마지막 헤드}(출력 1,000개
      $\to$ K개)뿐.
  \end{itemize}
\end{frame}
